\subsection{函子性}

设 \(f: \mathbf{C} \to \mathbf{C}'\) 是复形的态射。我们用

\[
\begin{array}{l} \mathbf {K} _ {\bullet} ^ {\mathrm{A}} (u, f): \mathbf {K} _ {\bullet} ^ {\mathrm{A}} (u, \mathrm{C}) \to \mathbf {K} _ {\bullet} ^ {\mathrm{A}} (u, \mathrm{C} ^ {\prime}), \\ \mathbf {K} _ {\mathrm{A}} ^ {\bullet} (u, f): \mathbf {K} _ {\mathrm{A}} ^ {\bullet} (u, \mathrm{C}) \to \mathbf {K} _ {\mathrm{A}} ^ {\bullet} (u, \mathrm{C} ^ {\prime}), \end{array}
\]

表示复形的态射 \(f \otimes 1_{\mathbf{K}(u)}\) 和 \(\operatorname{Homgr}_{\mathbf{A}}(1_{\mathbf{K}(u)}, f)\) 。

我们用 \(\mathrm{H}_{\bullet}^{\mathrm{A}}(u,f):\mathrm{H}_{\bullet}^{\mathrm{A}}(u,\mathbf{C})\to\mathrm{H}_{\bullet}^{\mathrm{A}}(u,\mathbf{C}^{\prime}),\quad\mathrm{H}_{\mathrm{A}}^{\bullet}(u,f):\mathrm{H}_{\mathrm{A}}^{\bullet}(u,\mathbf{C})\to\mathrm{H}_{\mathrm{A}}^{\bullet}(u,\mathbf{C}^{\prime}),\quad\mathrm{H}_{r}^{\mathrm{A}}(u,f):\mathrm{H}_{r}^{\mathrm{A}}(u,\mathbf{C})\to\mathrm{H}_{r}^{\mathrm{A}}(u,\mathbf{C}^{\prime}),\mathrm{H}_{\mathrm{A}}^{r}(u,f):\mathrm{H}_{\mathrm{A}}^{r}(u,\mathbf{C})\to\mathrm{H}_{\mathrm{A}}^{r}(u,\mathbf{C}^{\prime})\) 表示在同调中诱导的态射。映射 \(f\mapsto\mathbf{K}_{\bullet}^{\mathrm{A}}(u,f)\) 是线性的；如果 \(g:\mathbf{C}^{\prime}\to\mathbf{C}^{\prime\prime}\) 是复形的另一个态射，我们有 \(\mathbf{K}_{\bullet}^{\mathrm{A}}(u,g\circ f)=\mathbf{K}_{\bullet}^{\mathrm{A}}(u,g)\circ\mathbf{K}_{\bullet}^{\mathrm{A}}(u,f)\) ；类似地，对于 \(\mathbf{K}_{\mathrm{A}}^{\bullet},\mathrm{H}_{\bullet}^{\mathrm{A}},\mathrm{H}_{\mathrm{A}}^{\bullet},\mathrm{H}_{r}^{\mathrm{A}},\mathrm{H}_{\mathrm{A}}^{r}\) .

设 \(0 \to \mathbf{C}' \xrightarrow{f} \mathbf{C} \xrightarrow{g} \dot{\mathbf{C}}'' \to 0\) 是复形的一个正合列。

\begin{enumerate}
\def\labelenumi{\alph{enumi})}
\tightlist
\item
  假设 L 是平坦的；则 \(\Lambda(L)\) 是平坦的（X，第 15 页，推论）。序列
\end{enumerate}

\[
0 \rightarrow \mathbf {K} _ {\bullet} ^ {\mathrm{A}} (u, C ^ {\prime}) \xrightarrow {\mathbf {K} _ {\bullet} ^ {\mathrm{A}} (u , f)} \mathbf {K} _ {\bullet} ^ {\mathrm{A}} (u, C) \xrightarrow {\mathbf {K} _ {\bullet} ^ {\mathrm{A}} (u , g)} \mathbf {K} _ {\bullet} ^ {\mathrm{A}} (u, C ^ {\prime \prime}) \rightarrow 0
\]

于是它是正合的，并（X，第 30 页）导出一个同调正合列。

\[
\dots \to \mathrm{H} _ {n} ^ {\mathbf {A}} (u, \mathrm{C} ^ {\prime}) \xrightarrow {\mathrm{H} _ {n} (u , f)} \mathrm{H} _ {n} ^ {\mathbf {A}} (u, \mathrm{C}) \xrightarrow {\mathrm{H} _ {n} (u , g)} \mathrm{H} _ {n} ^ {\mathbf {A}} (u, \mathrm{C} ^ {\prime \prime}) \xrightarrow {\hat {\sigma} _ {n}} \mathrm{H} _ {n - 1} ^ {\mathbf {A}} (u, \mathrm{C} ^ {\prime}) \to \dots .
\]

\begin{enumerate}
\def\labelenumi{\alph{enumi})}
\setcounter{enumi}{1}
\tightlist
\item
  假设 L 是投射的；则 \(\Lambda(L)\) 是投射的。序列
\end{enumerate}

\[
0 \rightarrow \mathbf {K} _ {\mathrm{A}} ^ {*} (u, \mathrm{C} ^ {\prime}) \xrightarrow {\mathbf {K} _ {\mathrm{A}} ^ {*} (u , f)} \mathbf {K} _ {\mathrm{A}} ^ {*} (u, \mathrm{C}) \xrightarrow {\mathbf {K} _ {\mathrm{A}} ^ {*} (u , g)} \mathbf {K} _ {\mathrm{A}} ^ {*} (u, \mathrm{C} ^ {\prime \prime}) \rightarrow 0
\]

于是它是正合的，并导出一个同调正合列。

\[
\dots \rightarrow \mathrm{H} _ {\mathrm{A}} ^ {n} (u, \mathrm{C} ^ {\prime}) \xrightarrow {\mathrm{H} ^ {n} (u , f)} \mathrm{H} _ {\mathrm{A}} ^ {n} (u, \mathrm{C}) \xrightarrow {\mathrm{H} ^ {n} (u , g)} \mathrm{H} _ {\mathrm{A}} ^ {n} (u, \mathrm{C} ^ {\prime \prime}) \xrightarrow {\partial^ {n}} \mathrm{H} _ {\mathrm{A}} ^ {n + 1} (u, \mathrm{C} ^ {\prime}) \rightarrow \dots .
\]

设 \(\rho: A \to A'\) 是一个环同态， \(L'\) 为 \(A'\) -模 \(A' \otimes_A L\) ， \(u': L' \to A'\) 为线性形式 \(1 \otimes u\) 。典范双射同态（III，第 83 页，命题 8）

\[
\psi : \Lambda_ {A ^ {\prime}} (A ^ {\prime} \otimes_ {A} L) \rightarrow A ^ {\prime} \otimes_ {A} \Lambda_ {A} (L)
\]

是 A'-模复形的同构。我们推出：

\begin{enumerate}
\def\labelenumi{\arabic{enumi})}
\tightlist
\item
  对每个 A'-模复形 C'，一个 A-模复形的同构
\end{enumerate}

\[
\mathbf {K} _ {\bullet} ^ {\mathrm{A} ^ {\prime}} \left(u ^ {\prime}, \mathrm{C} ^ {\prime}\right)\rightarrow \mathbf {K} _ {\bullet} ^ {\mathrm{A}} \left(u, \mathrm{C} ^ {\prime}\right),
\]

\begin{enumerate}
\def\labelenumi{(\arabic{enumi})}
\setcounter{enumi}{8}
\tightlist
\item
\end{enumerate}

它由图表

\[
\mathrm{C} ^ {\prime} \otimes_ {\mathrm{A} ^ {\prime}} \left(\boldsymbol {\Lambda} _ {\mathrm{A} ^ {\prime}} \left(\mathrm{A} ^ {\prime} \otimes_ {\mathrm{A}} \mathrm{L}\right)\right) \xrightarrow {1 _ {\mathrm{C} ^ {\prime}} \otimes \psi} \mathrm{C} ^ {\prime} \otimes_ {\mathrm{A} ^ {\prime}} \mathrm{A} ^ {\prime} \otimes_ {\mathrm{A}} \boldsymbol {\Lambda} _ {\mathrm{A}} (\mathrm{L}) \xrightarrow {\varphi} \mathrm{C} ^ {\prime} \otimes_ {\mathrm{A}} \boldsymbol {\Lambda} _ {\mathrm{A}} (\mathrm{L})
\]

复合而成，其中 φ 是典范双射（III，第 85 页，命题 14）；

\begin{enumerate}
\def\labelenumi{\arabic{enumi})}
\setcounter{enumi}{1}
\tightlist
\item
  对每个 A-模复形 C，一个 A'-模复形的同构
\end{enumerate}

\[
\mathbf {K} _ {\bullet} ^ {\mathrm{A}} (u, \mathrm{A} ^ {\prime} \otimes_ {\mathrm{A}} \mathrm{C}) \rightarrow \mathrm{A} ^ {\prime} \otimes_ {\mathrm{A}} \mathbf {K} _ {\bullet} ^ {\mathrm{A}} (u, \mathrm{C}),
\]

由此得到A'-模的同态

\[
\mathrm{A} ^ {\prime} \otimes_ {\mathrm{A}} \mathrm{H} _ {n} ^ {\mathrm{A}} (u, \mathrm{C}) \rightarrow \mathrm{H} _ {n} ^ {\mathrm{A} ^ {\prime}} (u ^ {\prime}, \mathrm{A} ^ {\prime} \otimes_ {\mathrm{A}} \mathrm{C}),
\]

当 A' 在 A 上平坦时，这些同态是双射的（X，第 66 页，推论 2）。

设 L' 为一个 A-模， \(u': L' \to A\) 为一个线性形式， \(f: L \to L'\) 为一个 A-同态，使得 \(u' \circ f = u\) 。由 III，第 161 页，公式 (55)，可知同态 \(\symbf{\Lambda}(f): \symbf{\Lambda}(\mathrm{L}) \to \symbf{\Lambda}(\mathrm{L}')\) 满足 \(d_u \circ \symbf{\Lambda}(f) = \symbf{\Lambda}(f) \circ d_{u'}\) ，从而定义了复形的态射 \(\symbf{\Lambda}(u): \mathbf{K}^{\mathbf{A}}(u) \to \mathbf{K}^{\mathbf{A}}(u')\) 。若 C 是一个 A-复形，我们推出复形的态射

\[
1 _ {C} \otimes \boldsymbol {\Lambda} (u): \mathbf {K} ^ {\mathrm{A}} (u, C) \rightarrow \mathbf {K} ^ {\mathrm{A}} \left(u ^ {\prime}, C\right) \quad \text { and }\quad \operatorname{Homgr} \left(\boldsymbol {\Lambda} (u), 1 _ {C}\right): \mathbf {K} _ {A} ^ {*} \left(u ^ {\prime}, C\right)\rightarrow \mathbf {K} _ {A} ^ {*} (u, C)
\]

若f是双射，则所有这些态射都是同构。

